\section*{Chapitre \ref{ch:iv1:timed-test-8} --- Épreuve chronométrée 8}
\markboth{\protect\omTocMark{}{Réponses~: Épreuve chronométrée 8}}{\protect\omTocMark{}{Réponses~: Épreuve chronométrée 8}}

\begin{solution}{iq:iv1:timed-test-8:1}
La première est plus probable, avec $0.518$ contre $0.491$.

On passe au complémentaire\index{complémentaire}. La première vaut $1 - (\frac56)^4 = 1 - \frac{625}{1296} = 0.518$. La seconde vaut $1 - (\frac{35}{36})^{24} = 0.491$.

Ce qui rend la mauvaise réponse séduisante, c'est que les deux cas se ressemblent. Dans les deux cas, le nombre d'essais multiplié par la probabilité par essai vaut $\frac23$, puisque $4 \times \frac16 = 24 \times \frac{1}{36}$. Mais cette règle linéaire compte les chevauchements comme s'ils ne se produisaient pas, et les nombreux petits essais se chevauchent davantage que les quelques grands~: l'espérance du nombre de paires de succès vaut environ $0.21$ contre $0.17$. Un calcul exact confirme les deux résultats, et une simulation confirme ces valeurs.

À la limite d'un grand nombre de petits essais, la probabilité tend vers $1 - e^{-2/3} = 0.487$, si bien que la seconde en est déjà proche.

\iqwrong{Affirmer que les deux probabilités sont égales, car le nombre de succès attendu vaut $\frac23$ dans les deux cas. Une espérance n'est pas une probabilité.}
\end{solution}

\begin{solution}{iq:iv1:timed-test-8:2}
La probabilité vaut $\frac{9}{64} \approx 0.141$.

Par dénombrement direct~: les marches fermées comportent soit deux paires de pas opposés sur un même axe, soit une paire sur chaque axe. Deux paires sur un même axe~: $\binom42 = 6$ ordres pour chacun des 2 axes, soit 12. Une paire sur chaque axe~: $4! = 24$ ordres. Cela fait 36 marches sur $4^4 = 256$.

Une seconde voie fait tourner la grille de $45^\circ$\index{astuce de rotation}. Chaque pas déplace alors les deux nouvelles coordonnées de $\pm1$, indépendamment. Chaque coordonnée est une marche séparée, et chacune revient à son point de départ après 4 pas avec la probabilité $\binom42/16 = \frac38$. Le carré de ce nombre vaut $\frac{9}{64}$.

Une énumération des 256 marches le confirme. La probabilité est plus faible que le $\frac38$ d'une seule dimension, car il y a davantage de place pour errer.

\iqwrong{Ne compter que les marches qui s'éloignent deux fois puis reviennent tout droit. Les marches comme droite, haut, gauche, bas referment aussi la boucle, et l'ordre des pas compte.}
\end{solution}

\begin{solution}{iq:iv1:timed-test-8:3}
Le marché intègre une croissance de $4~\%$. À $5~\%$, le prix vaut environ 133.

La perpétuité croissante donne $P = \frac{D}{r - g}$\index{dividendes actualisés}. Donc $100 = \frac{4}{0.08 - g}$, et $g = 4\%$. Une somme numérique des dividendes actualisés donne le même prix de 100.

Pour $g = 5~\%$, $P = \frac{4}{0.03} = 133.3$. Un point de croissance supplémentaire fait monter le prix d'un tiers. C'est pourquoi les actifs à longue échéance sont si sensibles à la croissance et aux taux.

Vérification à la limite~: lorsque $g$ tend vers $r$, le prix tend vers l'infini, et la formule n'a plus de sens à partir de $r$ ou au-delà. Dans ce cas, la somme diverge, ce qui signifie que c'est l'hypothèse, et non le marché, qui est fausse.

\iqwrong{Considérer les 8~\% comme le seul rendement en dividende, et répondre donc une croissance nulle. Le rendement vaut 4~\%, et la croissance doit fournir le reste.}
\end{solution}

\begin{solution}{iq:iv1:timed-test-8:4}
La probabilité vaut $\frac{5}{36}$.

Une différence de 5 signifie qu'un 1 et un 6 apparaissent tous deux. On utilise l'inclusion--exclusion\index{inclusion-exclusion} sur le complémentaire.

Il y a 216 issues. Celles qui ne contiennent aucun 1 sont au nombre de $5^3 = 125$, et il en va de même pour celles sans 6. Celles qui n'en contiennent ni 1 ni 6 sont au nombre de $4^3 = 64$. Les issues favorables sont donc $216 - 125 - 125 + 64 = 30$. Et $\frac{30}{216} = \frac{5}{36}$.

Un dénombrement direct le vérifie~: le troisième dé vaut 1 ou 6, avec 3 ordres pour chacun, ce qui donne 6, ou une valeur de 2 à 5, avec 6 ordres pour chacune, ce qui donne 24. Le total vaut 30.

\iqwrong{Compter $3! \times 6$ pour le troisième dé dans tous les cas, ce qui donne 36. Lorsque le troisième dé répète un 1 ou un 6, il y a moins d'ordres.}
\end{solution}

\begin{solution}{iq:iv1:timed-test-8:5}
La probabilité vaut $\frac34$.

Sachant $X = x$, la probabilité que $Y < 1 - x$ vaut $1 - x$. Conditionner sur $X < \frac12$ rend $X$ uniforme sur $[0, \frac12]$, de moyenne $\frac14$. La réponse est donc $1 - \frac14 = \frac34$\index{conditionnement}.

Géométriquement, la région définie par la condition est un rectangle d'aire $\frac12$. La partie située sous la droite $x + y = 1$ est un trapèze d'aire $\frac38$. Le rapport vaut $\frac34$.

Une simulation confirme ce résultat. Vérification à la limite~: conditionner sur $X < a$ pour $a$ petit donne $1 - a/2$, qui tend vers 1, et pour $a = 1$ donne $\frac12$, la probabilité non conditionnelle.

\iqwrong{Donner $\frac12$, la valeur non conditionnelle. Savoir que $X$ est petit rend la somme plus petite, donc la probabilité doit augmenter.}
\end{solution}

\begin{solution}{iq:iv1:timed-test-8:6}
La probabilité vaut $\frac{7}{18}$.

Conditionner selon le type de la première tuile\index{conditionnement sur le premier tirage}. Une tuile double a une probabilité $\frac{7}{28} = \frac14$. Elle s'apparie avec les 6 autres tuiles portant ce nombre. Une tuile non double a une probabilité $\frac34$, et elle s'apparie avec 6 autres tuiles sur chacun de ses deux nombres, soit 12 tuiles.

Parmi les 27 tuiles restantes,
\[
\frac14 \cdot \frac{6}{27} + \frac34 \cdot \frac{12}{27} = \frac{6 + 36}{108} = \frac{7}{18}.
\]
Une énumération des $\binom{28}{2}$ paires donne 147 paires assorties sur 378, ce qui redonne la même valeur.

\iqwrong{Compter 12 correspondances pour chaque tuile, ce qui donne $\frac49$. Une tuile double n'a qu'un seul nombre à offrir.}
\end{solution}

\begin{solution}{iq:iv1:timed-test-8:7}
L'espérance du nombre de tirages vaut $e \approx 2.718$.

Soit $N$ le nombre de tirages. On a $N > n$ exactement lorsque les $n$ premiers tirages sont rangés dans l'ordre croissant, ce qui a pour probabilité $\frac{1}{n!}$, par symétrie entre les $n!$ ordres possibles\index{somme des queues pour l'espérance}. Donc
\[
\mathbb{E}N = \sum_{n \ge 0}P(N > n) = \sum_{n \ge 0}\frac{1}{n!} = e.
\]
Une simulation donne $2.718$.

Vérifions sur les petits cas~: $P(N = 2) = \frac12$, une descente dès le deuxième tirage, et $P(N = 3) = \frac12 - \frac16 = \frac13$. Ces deux seuls termes apportent 2 à la moyenne, si bien qu'un total juste inférieur à 3 est vraisemblable.

La réponse ne dépend pas de la loi, pourvu qu'elle soit continue, car seul compte l'ordre.

\iqwrong{Traiter chaque tirage ultérieur comme ayant une probabilité $\frac12$ d'être une descente, ce qui donne 3 par une loi géométrique. Une longue suite croissante rend la dernière valeur élevée, et donc la prochaine descente plus vraisemblable.}
\end{solution}

\begin{solution}{iq:iv1:timed-test-8:8}
Avec des parties équitables, la part de A vaut $\frac{11}{16}$. Avec $p = 0.6$, elle vaut environ $0.821$.

On joue toutes les parties qui pourraient être nécessaires, même si le match pourrait s'arrêter plus tôt. Au plus 4 parties suffisent à le décider\index{partie étendue}. A gagne exactement lorsqu'il remporte au moins 2 parties sur les 4, car B ne peut alors pas atteindre 3.

Avec des parties équitables, la probabilité vaut $1 - \frac{1 + 4}{16} = \frac{11}{16}$.

Avec $p = 0.6$, la probabilité que A remporte 0 ou 1 partie vaut $0.4^4 + 4 \times 0.6 \times 0.4^3 = 0.0256 + 0.1536 = 0.1792$. La part de A vaut donc $0.8208$. Une récurrence directe sur le score donne ces deux valeurs.

Vérification à la limite~: si A n'avait plus besoin d'aucune partie, sa part vaudrait 1, et si les deux joueurs avaient besoin du même nombre, elle vaudrait $\frac12$.

\iqwrong{Répartir selon le nombre de parties encore nécessaires, dans le rapport 3 à 2, ce qui donne $\frac35$. Cela ignore la façon dont les parties restantes peuvent tomber.}

\iqwhy{C'est le fondement de la valorisation par arbre~: la valeur aujourd'hui est la valeur pondérée de tout ce qui peut encore arriver.}
\end{solution}

\begin{solution}{iq:iv1:timed-test-8:9}
En dimension 3, le taux vaut $\frac{\pi}{6} \approx 0.524$, soit $\frac{6}{\pi} \approx 1.91$ tirages. En dimension 10, le taux vaut environ $0.0025$.

Le taux d'acceptation est le rapport des volumes\index{echantillonnage@échantillonnage par rejet}. La boule a pour volume $\frac43\pi$, et le cube, de côté 2, a pour volume 8. On obtient donc $\frac{\pi}{6}$. Le nombre de tirages suit une loi géométrique, de moyenne égale à l'inverse du taux.

En dimension 10, le cube a pour volume $2^{10} = 1024$. Et $\pi^5 \approx 306$, donc la boule a pour volume $\frac{306}{120} = 2.55$. Le rapport vaut $\frac{2.55}{1024} = 0.0025$, soit environ 400 tirages par point.

Une simulation confirme les deux valeurs. Le rejet depuis un cube échoue donc en grande dimension, car presque tout le cube se trouve dans ses coins. Il faut plutôt tirer un vecteur gaussien, le normaliser, puis redimensionner le rayon.

\iqwrong{S'attendre à ce que le taux se stabilise près de $\frac12$, parce que c'est le cas en faible dimension. Il décroît plus vite qu'exponentiellement avec la dimension.}
\end{solution}

\begin{solution}{iq:iv1:timed-test-8:10}
L'espérance vaut $\frac45$.

On considère un roi et les quatre as. Parmi ces cinq cartes, le roi a la même probabilité d'occuper chaque position\index{symetrie@symétrie}, donc il arrive en premier avec probabilité $\frac15$.

Chaque roi précède donc tous les as avec probabilité $\frac15$. En additionnant les variables indicatrices sur les 4 rois, on obtient $\frac45$.

Une simulation donne $0.80$. Une seconde voie~: les cartes qui ne sont pas des as et qui précèdent le premier as sont en moyenne au nombre de $\frac{48}{5}$, et les rois représentent 4 des 48 cartes, donc la moyenne vaut $\frac{48}{5} \cdot \frac{4}{48} = \frac45$.

On remarque que les 44 autres cartes du jeu n'interviennent pas, et que la première voie ne les a jamais utilisées.

\iqwrong{Prendre $\frac14$ par roi, ce qui donne 1. Chaque roi est en compétition avec 4 as, donc cinq cartes sont en jeu, et non quatre.}
\end{solution}

\begin{solution}{iq:iv1:timed-test-8:11}
L'espérance du nombre d'arrêts vaut $\frac{4651}{1296} \approx 3.59$, et la probabilité vaut $\frac{5}{54} \approx 0.093$.

Un étage n'est choisi par aucune des 5 personnes avec probabilité $(\frac56)^5 = \frac{3125}{7776}$. En sommant les variables indicatrices sur les 6 étages\index{variables indicatrices},
\[
\mathbb{E} = 6\left(1 - \frac{3125}{7776}\right) = \frac{4651}{1296}.
\]

Toutes différentes~: $6 \cdot 5 \cdot 4 \cdot 3 \cdot 2 = 720$ des $6^5 = 7776$ choix possibles, soit $\frac{5}{54}$.

Une énumération de tous les 7776 choix donne le même résultat. Vérification à la limite~: avec beaucoup d'étages, les arrêts tendent vers 5, un par personne. Avec un seul étage, l'expression vaut exactement 1.

\iqwrong{Répondre 5, une par personne, sans tenir compte des arrêts communs. Il faut raisonner étage par étage, et non personne par personne.}
\end{solution}

\begin{solution}{iq:iv1:timed-test-8:12}
Le $R^2$ avec les deux vaut $\frac13$, et avec un seul d'entre eux, il vaut $\frac14$.

En utilisant des variables centrées réduites\index{equations@équations normales}, les coefficients vérifient $b_1 + \frac12b_2 = \frac12$ et $\frac12b_1 + b_2 = \frac12$. Par symétrie, $b_1 = b_2 = \frac13$.

Et $R^2 = \sum b_ic_i = \frac13 \cdot \frac12 + \frac13 \cdot \frac12 = \frac13$.

Avec un seul, il vaut le carré de la corrélation, soit $\frac14$. Le second régresseur n'apporte donc que $\frac{1}{12}$. Une résolution numérique le confirme.

Vérifications à la limite~: si les régresseurs étaient non corrélés, le $R^2$ vaudrait $\frac14 + \frac14 = \frac12$, et s'ils étaient parfaitement corrélés, il resterait égal à $\frac14$. Notre valeur se situe entre les deux, comme il se doit.

\iqwrong{Additionner les carrés, ce qui donne $\frac12$. Cela n'est valable que pour des régresseurs non corrélés, car le recouvrement ne doit pas être compté deux fois.}
\end{solution}
